% ---------------------------------------------------------------------------
%  Validation study: the SDEM transport law against Fernandez de Labastida &
%  Yaroshchuk, Membranes 11 (2021) 272.  HAND-WRITTEN prose; every table is
%  \input{} from a GENERATED file (bin/curate/gen_sdem_validation.py), so no
%  number below was typed.  Written 2026-09-15 to be read by the paper's
%  authors.
% ---------------------------------------------------------------------------

\clearpage
\section*{Validation study: ion rejection in nanofiltration against Fern\'andez de Labastida \& Yaroshchuk (2021)}\addcontentsline{toc}{section}{Validation study: NF ion rejection (SDEM) against Fern\'andez de Labastida \& Yaroshchuk 2021}\markboth{Validation: SDEM}{Validation: SDEM}

This section compares one transport law of Choupo's membrane module, the
solution-diffusion-electromigration (SDEM) law, against one published data
set: M.~Fern\'andez de Labastida and A.~Yaroshchuk, \emph{Nanofiltration of
Multi-Ion Solutions: Quantitative Control of Concentration Polarization and
Interpretation by Solution-Diffusion-Electro-Migration Model}, Membranes 11
(2021) 272, doi:10.3390/membranes11040272, published under CC~BY~4.0
(reference~[3] in the list that closes this section; the model itself is
Yaroshchuk, Bruening and Lic\'on Bernal 2013~[1], and the analytical
trace-ion solution the paper fits with is Yaroshchuk and Bruening 2017~[2]).
It is written so that the authors of that paper can check it.  Every number in the
tables is produced by the engine or read from the published figures by a
script; none is typed.  The section ends with the commands that regenerate
it and with the points on which we would value the authors' reading.

\subsection*{The model as Choupo implements it}

The active layer is a slab of dimensionless thickness $x \in [0,1]$.  Each
ion $i$ has a membrane permeance $P_i$ (m/s), which already contains its
partition coefficient (the solution-diffusion definition $P = DK/\delta$), a
charge $z_i$, and a \emph{virtual} concentration $c_i(x)$: the concentration
of the solution that would be in local equilibrium with the membrane at $x$.
The flux of each ion is Nernst--Planck on the virtual solution,
\[
  j_i \;=\; -P_i\left(\frac{\mathrm{d}c_i}{\mathrm{d}x} + z_i\,c_i\,\frac{\mathrm{d}\psi}{\mathrm{d}x}\right),
  \qquad \psi = \frac{F\,\varphi}{RT},
\]
the virtual solution is electroneutral at every $x$, $\sum_i z_i c_i(x) = 0$,
and no current flows, $\sum_i z_i j_i = 0$.  The last two conditions fix the
one electric field the ions share,
\[
  \frac{\mathrm{d}\psi}{\mathrm{d}x} \;=\; -\,\frac{\sum_i z_i\, j_i / P_i}{\sum_i z_i^2\, c_i},
\]
so there is exactly one state variable beyond the uncoupled solution-diffusion
law, the potential drop $\Delta\psi$ across the layer, and no new parameter.
The boundary conditions are the wall concentrations $c_i(0) = c_{m,i}$
(from the film model, per ion) and the permeate condition $j_i = J_v\,c_{p,i}$
with $c_i(1) = c_{p,i}$, where $J_v$ is the transmembrane volume flux.  This
is the formulation of Yaroshchuk, Bruening and Lic\'on Bernal~[1], and it has
a property the engine checks rather than
assumes: for a single salt it reduces \emph{exactly} to
\[
  R_s \;=\; \frac{J_v}{J_v + P_s}, \qquad
  P_s \;=\; \frac{(z_+ + |z_-|)\,P_+\,P_-}{z_+ P_+ + |z_-|\,P_-},
\]
the ambipolar permeance, which is what makes per-ion permeances calibratable
from single-salt tests.

What the implementation does \emph{not} contain, so that the comparison is
read for what it is: no fixed membrane charge (no Donnan partition, no
dielectric exclusion --- those effects are inside the permeances the authors
fitted); ideal dilute virtual solutions; constant permeances; and a
concentration-polarisation film that is computed per ion \emph{without} a
field (the engine announces this every run).  All runs below use a film
coefficient of 1~m/s, so the wall equals the bulk to better than $0.01\,\%$
and the rejections are intrinsic --- which is what the paper's CP-corrected
rejections are.  Numerically, each station solves a Newton problem in the
permeate composition whose residual is a fourth-order Runge--Kutta shot of
the profile \emph{from the permeate face back to the wall} (320 steps); the
osmotic back-pressure closes an outer Newton in $J_v$.  The source is
\file{src/unitOperations/membrane/transport/SDEM.cpp}.

\subsection*{What was taken from the paper}

\textbf{The permeances.}  Table~1 of the paper gives, for each dominant
salt, the membrane permeance to the salt $P_s$ fitted with Eq.~(1) and the
ion permeances fitted with the SDEM trace solution (its Eq.~(3), the
analytical solution of~[2]), all in
$\mu$m/s.  Choupo uses the ion values as they stand.  Where the paper gives
a \emph{bound} or a \emph{range} the value used is stated in the case
record: Na$^+$ under NaCl ``20--113'' is taken at 113, the end that
reproduces $P_s = 7.5$ through the ambipolar identity; NH$_4^+$ ``$>60$''
(NaCl), ``$>700$'' (MgCl$_2$), ``$>100$'' (Na$_2$SO$_4$), Na$^+$
``$>250$'' (MgCl$_2$) and ``$>0.8$'' (Na$_2$SO$_4$) are taken at the
bound.  The consequences of the bounds are measured below rather than
assumed away.

\textbf{The feeds.}  Dominant salt $0.01$~mol/L, the two trace salts
NaNO$_3$ and NH$_4$Cl at $2\times10^{-4}$~mol/L each, $20\,^\circ$C, on a
flat NF270 disc in the rotating-disk cell, transmembrane pressure 2--14~bar.

\textbf{The figures.}  Figures~1--4 plot, for each dominant salt, the
reciprocal intrinsic transmission $f = 1/(1-R)$ of the salt (panel~a) and of
the trace ions (panel~b) against $J_v$.  The \emph{symbols} are the
measurements after the authors' CP correction; the \emph{lines} are the
authors' fits --- Eq.~(1) for the salt, the analytical trace-limit solution
Eq.~(3) for a trace --- with the Table~1 permeances.  The publisher's
full-resolution panels were digitised by
\file{bin/curate/digitise\_fdl2021.py}: the axes are calibrated from the
axis lines and verified against the tick labels (they map back to
$0, 10, 20, \dots$ within a few hundredths of a $\mu$m/s), markers are
separated from lines by a morphological opening and located at the centre
of their bounding box, the drawn line is read under each marker, and the
three arrows of Fig.~3b are blanked as annotations.  The reading error is
about $\pm 0.07\,\mu$m/s in $J_v$ and $\pm 0.3\,\%$ of the axis span in
$f$; a marker partly hidden under another series' symbol carries a larger
error and is marked $\dagger$ in the tables.  The 106 points live, with
their reading errors, in the NaCl witness case:\par\noindent
{\footnotesize\path{tutorials/steady/membranes/membrane12_sdem_nf270_nacl_traces/constant/experimental/fdl2021_figures.csv}}\par\noindent
The images are not redistributed.  One reading
we ask the authors to confirm: the right-hand axis of Fig.~3b (Cl$^-$ and
NO$_3^-$ under Na$_2$SO$_4$) is labelled 0 to 8 but carries nine equal tick
intervals; it was read as 0 to 18.

\subsection*{How Choupo was driven}

For each digitised point the engine is run on a one-element NF270 case built
from the shipped NaCl witness (\file{membrane12\_sdem\_nf270\_nacl\_traces})
with the salt's Table~1 row as the membrane record, and the feed pressure is
adjusted until the engine's $J_v$ equals the figure's to $0.5\,\%$.  Two
feeds are used for every salt.  \textbf{Trace limit}: the trace salts at
$2\times10^{-7}$~mol/L, one thousandth of the paper's, which is the limit
the analytical Eq.~(3) assumes (a trace that does not feed back on the
field).  \textbf{Paper's feed}: the trace salts at $2\times10^{-4}$~mol/L,
the concentration actually used in the experiments.  The comparison that
tests the \emph{implementation} is therefore Choupo-in-the-trace-limit
against the authors' \emph{lines} (same model, same parameters, one answer);
the comparison that tests the \emph{model} is Choupo-at-the-paper's-feed
against the \emph{symbols}, and it inherits the authors' own fit residual.

\subsection*{Table 1 against the ambipolar identity}

The first check needs no run: the ion permeances of Table~1 must reproduce
the salt permeances of the same table through the ambipolar identity.  They
do, to the table's rounding.

% GENERATED by bin/curate/gen_sdem_validation.py from the engine and
% tutorials/steady/membranes/membrane12_sdem_nf270_nacl_traces/constant/experimental/fdl2021_figures.csv
% -- do not edit; regenerate.  Checked for drift by the same script with --check.
\begin{center}\footnotesize
\begin{tabular}{@{}llrlrrrr@{}}\toprule
salt & cation & $P_+$ & anion & $P_-$ & $P_s$ (Table 1) & $P_s$ ambipolar & dev.\ (\%) \\ \midrule
NaCl & Na$^+$ & 113 & Cl$^-$ & 3.9 & 7.5 & 7.54 & +0.5 \\
MgCl$_2$ & Mg$^{2+}$ & 3.7 & Cl$^-$ & 9.3 & 6.1 & 6.18 & +1.3 \\
Na$_2$SO$_4$ & Na$^+$ & 0.8 & SO$_4^{2-}$ & 0.06 & 0.16 & 0.16 & -2.2 \\
MgSO$_4$ & Mg$^{2+}$ & 0.9 & SO$_4^{2-}$ & 0.13 & 0.23 & 0.23 & -1.2 \\
\bottomrule
\end{tabular}
\end{center}

\noindent Permeances in $\mu$m/s.  This is arm (b) of the gate
\file{bin/curate/check\_sdem.py}, recomputed in Python independently of the
engine.

\subsection*{Dominant NaCl (Figure 1)}

% GENERATED by bin/curate/gen_sdem_validation.py from the engine and
% tutorials/steady/membranes/membrane12_sdem_nf270_nacl_traces/constant/experimental/fdl2021_figures.csv
% -- do not edit; regenerate.  Checked for drift by the same script with --check.
\begin{footnotesize}\setlength{\tabcolsep}{4pt}
\begin{longtable}{@{}lrrrrrr@{}}
\caption{Dominant NaCl, Figure 1 of the paper.  $f = 1/(1-R)$; $J_v$ in $\mu$m/s.  ``measured'' and ``paper's fit'' are read from the figure ($\dagger$: marker partly hidden, larger reading error; --: the drawn line stops before the point); ``trace limit'' and ``paper's feed'' are Choupo at the same flux with the trace salts at $2\times10^{-7}$ and $2\times10^{-4}$~mol/L; the last column is Choupo (trace limit) against the paper's fit.}\\
\toprule
ion & $J_v$ & $f$ measured & $f$ paper's fit & $f$ Choupo, trace limit & $f$ Choupo, paper's feed & dev.\ (\%) \\ \midrule
\endfirsthead
\toprule
ion & $J_v$ & $f$ measured & $f$ paper's fit & $f$ Choupo, trace limit & $f$ Choupo, paper's feed & dev.\ (\%) \\ \midrule
\endhead
\bottomrule
\endfoot
Na$^+$ (dominant) & 8.2 & 1.758 & 2.121 & 2.092 & 2.060 & -1.4 \\
 & 12.3 & 2.145 & 2.666 & 2.636 & 2.581 & -1.1 \\
 & 20.2 & 3.726 & 3.706 & 3.677 & 3.573 & -0.8 \\
 & 27.7 & 4.870 & 4.715 & 4.668 & 4.516 & -1.0 \\
 & 36.2 & 5.980 & 5.841 & 5.772 & 5.567 & -1.2 \\
 & 45.1 & 7.049 & 7.032 & 6.938 & 6.676 & -1.3 \\
 & 55.8 & 6.275 & -- & 8.336 & 8.007 & -- \\
\addlinespace[2pt]
Cl$^-$ (dominant) & 8.2 & 1.758 & 2.121 & 2.092 & 2.109 & -1.4 \\
 & 12.3 & 2.145 & 2.666 & 2.636 & 2.660 & -1.1 \\
 & 20.2 & 3.726 & 3.706 & 3.677 & 3.710 & -0.8 \\
 & 27.7 & 4.870 & 4.715 & 4.668 & 4.708 & -1.0 \\
 & 36.2 & 5.980 & 5.841 & 5.772 & 5.818 & -1.2 \\
 & 45.1 & 7.049 & 7.032 & 6.938 & 6.988 & -1.3 \\
 & 55.8 & 6.275 & -- & 8.336 & 8.391 & -- \\
\addlinespace[2pt]
NH$_4^+$ (trace) & 8.3 & 2.041 & 2.020 & 2.185 & 2.152 & +8.2 \\
 & 12.3 & 2.396 & 2.481 & 2.775 & 2.719 & +11.9 \\
 & 20.1 & 3.461 & 3.388 & 3.945 & 3.838 & +16.5 \\
 & 27.8 & 4.361 & 4.260 & 5.076 & 4.919 & +19.2 \\
 & 36.2 & 4.997 & 5.220 & 6.343 & 6.129 & +21.5 \\
 & 45.1 & 5.533 & 6.219 & 7.687 & 7.414 & +23.6 \\
 & 55.9 & 5.153 & -- & 9.315 & 8.970 & -- \\
\addlinespace[2pt]
NO$_3^-$ (trace) & 8.2 & 0.916 & 1.036 & 0.952 & 0.962 & -8.2 \\
 & 12.4 & 0.984 & 1.118 & 1.036 & 1.049 & -7.3 \\
 & 20.2 & 1.315 & 1.321 & 1.253 & 1.269 & -5.1 \\
 & 27.7 & 1.526 & 1.547 & 1.491 & 1.508 & -3.6 \\
 & 36.2 & 1.913 & 1.817 & 1.775 & 1.793 & -2.3 \\
 & 45.1 & 2.143 & 2.113 & 2.083 & 2.102 & -1.4 \\
 & 55.9 & 1.963 & 2.477 & 2.461 & 2.481 & -0.6 \\
\addlinespace[2pt]
\end{longtable}
\end{footnotesize}


For the salt, Choupo's dominant anion follows the paper's line to within
$1.5\,\%$ over the whole flux range (Cl$^-$ at $2\times10^{-4}$~M: the salt
rejection is set by the ambipolar permeance, and the traces perturb it
little here).  For NO$_3^-$ the trace-limit engine sits within $8\,\%$ of
the line, closest where the line is flattest; at the two lowest fluxes the
engine gives a slightly \emph{negative} rejection where the drawn line is
slightly positive --- and the measured symbols there are negative, as the
paper's text says (``lightly negative values at small transmembrane
fluxes'').  NH$_4^+$ is the case where the permeance is a bound: at
$60\,\mu$m/s the engine's $f$ is 8--24\,\% above the line, and the
sensitivity table further down shows the line is met only near
$1000\,\mu$m/s; the authors report that their fit ``became insensitive to the
value of ionic permeance as this increased'', which is consistent with a
fit that had already saturated.

\subsection*{Dominant MgCl$_2$ (Figure 2)}

% GENERATED by bin/curate/gen_sdem_validation.py from the engine and
% tutorials/steady/membranes/membrane12_sdem_nf270_nacl_traces/constant/experimental/fdl2021_figures.csv
% -- do not edit; regenerate.  Checked for drift by the same script with --check.
\begin{footnotesize}\setlength{\tabcolsep}{4pt}
\begin{longtable}{@{}lrrrrrr@{}}
\caption{Dominant MgCl$_2$, Figure 2 of the paper.  $f = 1/(1-R)$; $J_v$ in $\mu$m/s.  ``measured'' and ``paper's fit'' are read from the figure ($\dagger$: marker partly hidden, larger reading error; --: the drawn line stops before the point); ``trace limit'' and ``paper's feed'' are Choupo at the same flux with the trace salts at $2\times10^{-7}$ and $2\times10^{-4}$~mol/L; the last column is Choupo (trace limit) against the paper's fit.}\\
\toprule
ion & $J_v$ & $f$ measured & $f$ paper's fit & $f$ Choupo, trace limit & $f$ Choupo, paper's feed & dev.\ (\%) \\ \midrule
\endfirsthead
\toprule
ion & $J_v$ & $f$ measured & $f$ paper's fit & $f$ Choupo, trace limit & $f$ Choupo, paper's feed & dev.\ (\%) \\ \midrule
\endhead
\bottomrule
\endfoot
Mg$^{2+}$ (dominant) & 8.7 & 2.237 & 2.068 & 2.416 & 2.471 & +16.8 \\
 & 13.3 & 2.757 & 2.603 & 3.156 & 3.267 & +21.2 \\
 & 19.8 & 3.230 & 3.292 & 4.189 & 4.407 & +27.2 \\
 & 28.4 & 4.040 & 4.143 & 5.562 & 5.975 & +34.3 \\
 & 36.5 & 4.757 & 4.874 & 6.865 & 7.512 & +40.9 \\
 & 44.4 & 5.665 & 5.520 & 8.129 & 9.043 & +47.3 \\
 & 52.5 & 6.125 & 6.120 & 9.422 & 10.6 & +53.9 \\
\addlinespace[2pt]
Cl$^-$ (dominant) & 8.7 & 2.237 & 2.068 & 2.416 & 2.376 & +16.8 \\
 & 13.3 & 2.757 & 2.603 & 3.156 & 3.072 & +21.2 \\
 & 19.8 & 3.230 & 3.292 & 4.188 & 4.013 & +27.2 \\
 & 28.4 & 4.040 & 4.143 & 5.562 & 5.220 & +34.2 \\
 & 36.5 & 4.757 & 4.874 & 6.864 & 6.325 & +40.8 \\
 & 44.4 & 5.665 & 5.520 & 8.127 & 7.365 & +47.2 \\
 & 52.5 & 6.125 & 6.120 & 9.419 & 8.403 & +53.9 \\
\addlinespace[2pt]
NH$_4^+$ (trace) & 8.4 & 0.738$^\dagger$ & 0.798 & 0.760 & 0.776 & -4.7 \\
 & 13.3 & 0.684 & 0.694 & 0.698 & 0.725 & +0.5 \\
 & 19.8 & 0.604 & 0.623 & 0.643 & 0.686 & +3.2 \\
 & 28.4 & 0.548 & 0.567 & 0.597 & 0.662 & +5.3 \\
 & 36.5 & 0.532 & 0.529 & 0.568 & 0.654 & +7.4 \\
 & 44.5 & 0.485 & 0.500 & 0.549 & 0.654 & +9.6 \\
 & 52.5 & 0.476 & 0.478 & 0.534 & 0.658 & +11.7 \\
\addlinespace[2pt]
Na$^+$ (trace) & 19.8 & 0.724 & 0.695 & 0.685 & 0.730 & -1.4 \\
 & 28.4 & 0.679 & 0.668 & 0.656 & 0.724 & -1.8 \\
 & 36.5 & 0.671 & 0.657 & 0.644 & 0.733 & -2.0 \\
 & 44.4 & 0.640 & 0.656 & 0.639 & 0.750 & -2.5 \\
 & 52.5 & 0.651 & 0.661 & 0.641 & 0.772 & -3.0 \\
\addlinespace[2pt]
NO$_3^-$ (trace) & 8.7 & 1.603 & 1.891 & 1.933 & 1.899 & +2.2 \\
 & 13.3 & 1.856 & 2.309 & 2.396 & 2.325 & +3.8 \\
 & 19.7 & 2.830 & 2.887 & 3.029 & 2.886 & +4.9 \\
 & 28.4 & 3.421 & 3.650 & 3.859 & 3.589 & +5.7 \\
 & 36.5 & 4.243 & 4.354 & 4.636 & 4.220 & +6.5 \\
 & 44.4 & 5.063 & 5.025 & 5.381 & 4.805 & +7.1 \\
 & 52.5 & 6.093 & 5.674 & 6.136 & 5.382 & +8.1 \\
\addlinespace[2pt]
\end{longtable}
\end{footnotesize}


The three trace lines are reproduced in the trace limit within 1--8\,\%
(NO$_3^-$ $+2$ to $+8\,\%$, Na$^+$ $-1$ to $-3\,\%$, NH$_4^+$ $-5$ to
$+12\,\%$), with the negative rejections of both trace cations coming out
of the same permeances and no parameter for the effect.  The salt panel is
different, and it is a point for the authors: the drawn line of Fig.~2a
reads $f_s = 6.1$ at $52.5\,\mu$m/s and $5.5$ at $44.4\,\mu$m/s, whereas
Eq.~(1) with the Table~1 value $P_s = 6.1\,\mu$m/s gives $9.6$ and $8.3$
(off the top of the plotted axis); the measured points follow the drawn
line, and read on their own they imply an apparent salt permeance rising
from about 7 to 10~$\mu$m/s across the range.  Choupo, which solves the
same equations as Eq.~(1) for the dominant ions, agrees with Eq.~(1) and
the ambipolar identity (both dominant ions, up to the trace coupling
discussed below) and therefore \emph{not} with the drawn line.  We have
not tried to reconcile the two and would welcome the authors' reading of
which curve Fig.~2a shows.

\subsection*{Dominant Na$_2$SO$_4$ (Figure 3)}

% GENERATED by bin/curate/gen_sdem_validation.py from the engine and
% tutorials/steady/membranes/membrane12_sdem_nf270_nacl_traces/constant/experimental/fdl2021_figures.csv
% -- do not edit; regenerate.  Checked for drift by the same script with --check.
\begin{footnotesize}\setlength{\tabcolsep}{4pt}
\begin{longtable}{@{}lrrrrrr@{}}
\caption{Dominant Na$_2$SO$_4$, Figure 3 of the paper.  $f = 1/(1-R)$; $J_v$ in $\mu$m/s.  ``measured'' and ``paper's fit'' are read from the figure ($\dagger$: marker partly hidden, larger reading error; --: the drawn line stops before the point); ``trace limit'' and ``paper's feed'' are Choupo at the same flux with the trace salts at $2\times10^{-7}$ and $2\times10^{-4}$~mol/L; the last column is Choupo (trace limit) against the paper's fit.}\\
\toprule
ion & $J_v$ & $f$ measured & $f$ paper's fit & $f$ Choupo, trace limit & $f$ Choupo, paper's feed & dev.\ (\%) \\ \midrule
\endfirsthead
\toprule
ion & $J_v$ & $f$ measured & $f$ paper's fit & $f$ Choupo, trace limit & $f$ Choupo, paper's feed & dev.\ (\%) \\ \midrule
\endhead
\bottomrule
\endfoot
Na$^+$ (dominant) & 6.4 & 52.2 & 39.8 & 42.0 & 22.8 & +5.5 \\
 & 10.5 & 77.2 & 65.1 & 67.6 & 34.9 & +3.9 \\
 & 18.9 & 122.8 & 117.0 & 121.2 & 60.8 & +3.6 \\
 & 27.2 & 160.2 & 168.4 & 174.0 & 86.6 & +3.4 \\
 & 33.6 & 214.7 & 207.3 & 214.3 & 106.4 & +3.3 \\
 & 41.8 & 261.7 & 257.3 & 265.5 & 131.8 & +3.2 \\
 & 49.5 & 285.0 & 304.7 & 314.0 & 156.0 & +3.0 \\
\addlinespace[2pt]
SO$_4^{2-}$ (dominant) & 6.4 & 52.2 & 39.8 & 42.0 & 49.6 & +5.7 \\
 & 10.5 & 77.2 & 65.1 & 67.7 & 80.3 & +4.0 \\
 & 18.9 & 122.8 & 117.0 & 121.4 & 144.3 & +3.7 \\
 & 27.2 & 160.2 & 168.4 & 174.3 & 207.0 & +3.5 \\
 & 33.6 & 214.7 & 207.3 & 214.6 & 254.5 & +3.5 \\
 & 41.8 & 261.7 & 257.3 & 265.9 & 315.0 & +3.3 \\
 & 49.5 & 285.0 & 304.7 & 314.4 & 372.0 & +3.2 \\
\addlinespace[2pt]
NH$_4^+$ (trace) & 6.5 & 44.8 & 42.5 & 20.4 & 7.475 & -51.9 \\
 & 10.5 & 50.9 & 48.9 & 29.9 & 9.290 & -38.7 \\
 & 18.9 & 63.8 & 61.1 & 48.0 & 12.7 & -21.5 \\
 & 27.2 & 71.2 & 72.1 & 64.0 & 15.6 & -11.2 \\
 & 33.6 & 77.8 & 80.3 & 76.0 & 17.7 & -5.4 \\
 & 41.7 & 89.0 & 90.0 & 90.2 & 20.1 & +0.2 \\
 & 49.5 & 89.5 & 99.1 & 103.4 & 22.4 & +4.3 \\
\addlinespace[2pt]
Cl$^-$ (trace) & 6.5 & 1.337 & 1.113 & 1.127 & 1.349 & +1.2 \\
 & 10.5 & 1.972 & 1.663 & 1.761 & 2.061 & +5.9 \\
 & 18.9 & 3.384 & 2.896 & 3.119 & 3.587 & +7.7 \\
 & 27.3 & 4.429 & 4.130 & 4.467 & 5.098 & +8.1 \\
 & 33.6 & 6.735 & 5.094 & 5.490 & 6.244 & +7.8 \\
 & 41.7 & 7.354 & 6.274 & 6.796 & 7.705 & +8.3 \\
 & 49.5 & 7.036 & 7.422 & 8.040 & 9.096 & +8.3 \\
\addlinespace[2pt]
NO$_3^-$ (trace) & 6.4 & 0.501 & 0.514 & 0.423 & 0.557 & -17.7 \\
 & 10.5 & 0.719 & 0.692 & 0.634 & 0.788 & -8.3 \\
 & 18.9 & 1.137 & 1.149 & 1.098 & 1.299 & -4.5 \\
 & 27.2 & 1.437 & 1.613 & 1.561 & 1.812 & -3.2 \\
 & 33.6 & 2.089 & 1.965 & 1.918 & 2.209 & -2.4 \\
 & 41.7 & 2.524 & 2.418 & 2.371 & 2.712 & -2.0 \\
 & 49.5 & 2.808 & 2.842 & 2.800 & 3.190 & -1.5 \\
\addlinespace[2pt]
\end{longtable}
\end{footnotesize}


Here the traces are, in the authors' words, not authentic traces: sulfate
passes at $0.06\,\mu$m/s, the trace anions at 3.4 and 9.8, so the anion
current is carried by the traces.  In the trace limit Choupo reproduces
the Cl$^-$ line within $+1$ to $+8\,\%$ and the NO$_3^-$ line within
$-2$ to $-8\,\%$ from $10\,\mu$m/s upward ($-18\,\%$ at the lowest point,
where $f \approx 0.5$ and the drawn line is thick) (the reversal from negative to
positive rejection of NO$_3^-$ falls at about $16\,\mu$m/s in both); at the
paper's own feed the engine's $f$ for the trace anions is 8--24\,\% above
the lines and its $f$ for the dominant Na$^+$ well below that of
SO$_4^{2-}$ --- the effect the paper's text predicts (``the reciprocal
transmission of the more permeable ion (such as Na$^+$ in Na$_2$SO$_4$)
could be basically lower than that for SO$_4^{2-}$'').  NH$_4^+$ is not a
test of the implementation: its own permeance is a bound and, as the
sensitivity table shows, its rejection is set mostly by the Na$^+$
permeance, which the paper calls ``just orientative''.

\subsection*{Dominant MgSO$_4$ (Figure 4)}

% GENERATED by bin/curate/gen_sdem_validation.py from the engine and
% tutorials/steady/membranes/membrane12_sdem_nf270_nacl_traces/constant/experimental/fdl2021_figures.csv
% -- do not edit; regenerate.  Checked for drift by the same script with --check.
\begin{footnotesize}\setlength{\tabcolsep}{4pt}
\begin{longtable}{@{}lrrrrrr@{}}
\caption{Dominant MgSO$_4$, Figure 4 of the paper.  $f = 1/(1-R)$; $J_v$ in $\mu$m/s.  ``measured'' and ``paper's fit'' are read from the figure ($\dagger$: marker partly hidden, larger reading error; --: the drawn line stops before the point); ``trace limit'' and ``paper's feed'' are Choupo at the same flux with the trace salts at $2\times10^{-7}$ and $2\times10^{-4}$~mol/L; the last column is Choupo (trace limit) against the paper's fit.}\\
\toprule
ion & $J_v$ & $f$ measured & $f$ paper's fit & $f$ Choupo, trace limit & $f$ Choupo, paper's feed & dev.\ (\%) \\ \midrule
\endfirsthead
\toprule
ion & $J_v$ & $f$ measured & $f$ paper's fit & $f$ Choupo, trace limit & $f$ Choupo, paper's feed & dev.\ (\%) \\ \midrule
\endhead
\bottomrule
\endfoot
Mg$^{2+}$ (dominant) & 7.2 & 42.9 & 32.6 & 32.8 & 20.1 & +0.5 \\
 & 11.6 & 53.3 & 52.9 & 51.8 & 28.6 & -2.1 \\
 & 19.2 & 92.6 & 87.0 & 85.0 & 43.1 & -2.3 \\
 & 26.9 & 120.5 & 121.6 & 118.5 & 58.0 & -2.6 \\
 & 34.7 & 152.3 & 157.0 & 152.6 & 73.6 & -2.8 \\
 & 43.1 & 184.3 & 194.8 & 189.0 & 90.7 & -3.0 \\
 & 53.8 & 200.5 & 242.2 & 235.0 & 113.0 & -3.0 \\
\addlinespace[2pt]
SO$_4^{2-}$ (dominant) & 7.2 & 42.9 & 32.6 & 32.8 & 36.9 & +0.6 \\
 & 11.6 & 53.3 & 52.9 & 51.9 & 59.9 & -1.9 \\
 & 19.2 & 92.6 & 87.0 & 85.2 & 100.2 & -2.1 \\
 & 26.9 & 120.5 & 121.6 & 118.7 & 141.0 & -2.4 \\
 & 34.7 & 152.3 & 157.0 & 153.0 & 182.2 & -2.6 \\
 & 43.1 & 184.3 & 194.8 & 189.5 & 225.6 & -2.7 \\
 & 53.8 & 200.5 & 242.2 & 235.6 & 279.8 & -2.7 \\
\addlinespace[2pt]
NH$_4^+$ (trace) & 7.2 & 3.838$^\dagger$ & -- & 3.804 & 2.503 & -- \\
 & 19.2 & 5.212 & 5.429 & 5.618 & 2.838 & +3.5 \\
 & 53.8 & 7.651 & 8.416 & 8.686 & 3.825 & +3.2 \\
\addlinespace[2pt]
Na$^+$ (trace) & 7.2 & 3.427 & 3.643 & 3.815 & 2.506 & +4.7 \\
 & 11.4 & 4.554 & 4.345 & 4.563 & 2.634 & +5.0 \\
 & 19.2 & 5.698 & 5.358 & 5.613 & 2.838 & +4.8 \\
 & 26.9 & 6.080 & 6.161 & 6.459 & 3.049 & +4.8 \\
 & 34.8 & 6.838$^\dagger$ & 6.898 & 7.194 & 3.271 & +4.3 \\
 & 43.1 & 7.408 & 7.543 & 7.891 & 3.515 & +4.6 \\
 & 53.8 & 7.073 & 8.316 & 8.694 & 3.832 & +4.5 \\
\addlinespace[2pt]
Cl$^-$ (trace) & 7.2 & 0.712 & 0.660 & 0.615 & 0.788 & -6.8 \\
 & 11.6 & 0.825 & 0.817 & 0.774 & 1.003 & -5.3 \\
 & 19.2 & 1.215 & 1.127 & 1.088 & 1.386 & -3.5 \\
 & 26.9 & 1.492 & 1.459 & 1.420 & 1.771 & -2.7 \\
 & 34.7 & 1.827 & 1.807 & 1.767 & 2.163 & -2.2 \\
 & 43.1 & 2.137 & 2.185 & 2.141 & 2.577 & -2.0 \\
 & 53.8 & 2.313 & 2.665 & 2.618 & 3.099 & -1.8 \\
\addlinespace[2pt]
NO$_3^-$ (trace) & 7.2 & 0.419 & 0.440 & 0.401 & 0.558 & -8.9 \\
 & 11.6 & 0.478 & 0.471 & 0.434 & 0.631 & -7.9 \\
 & 19.2 & 0.587 & 0.559 & 0.528 & 0.767 & -5.7 \\
 & 26.9 & 0.687 & 0.668 & 0.640 & 0.904 & -4.2 \\
 & 34.7 & 0.788 & 0.785 & 0.761 & 1.042 & -3.1 \\
 & 43.1 & 0.913 & 0.916 & 0.894 & 1.189 & -2.4 \\
 & 53.8 & 0.955 & 1.096 & 1.067 & 1.374 & -2.7 \\
\addlinespace[2pt]
\end{longtable}
\end{footnotesize}


This is the case where the paper's fit determines every permeance, and in
the trace limit the engine reproduces all four trace lines within 2--9\,\%
(Na$^+$ $+4$ to $+5\,\%$; Cl$^-$ $-2$ to $-7\,\%$; NO$_3^-$ $-2$ to
$-9\,\%$; NH$_4^+$, whose line lies on Na$^+$'s, $+3\,\%$ on the points
visible under the Na$^+$ diamonds).  At the paper's own feed the engine
departs from the lines by 20--55\,\%, all in one direction: the trace
cations are rejected less and the trace anions less negatively.  Both are
the signature of a \emph{weaker} field, and the next subsection measures
why.

\subsection*{A trace is not a trace when the dominant anion is a hundred times slower}

The authors' Eq.~(3), from~[2], is the analytical solution for a trace ion in the
field set by the dominant salt alone.  At $2\times10^{-4}$~mol/L against
$0.01$~mol/L that assumption holds within $3\,\%$ for NaCl and within
$15\,\%$ for MgCl$_2$; under a sulfate it fails, because $2\times10^{-4}$~mol/L of NO$_3^-$ at
$42\,\mu$m/s carries more anion current than $0.01$~mol/L of SO$_4^{2-}$
at $0.13\,\mu$m/s.  The traces then relieve the field, and every rejection
moves toward its uncoupled value.  The table gives, for each salt at its
fifth measured flux, the ratio of the engine's $f$ at the paper's feed to
its $f$ in the trace limit, per trace ion.

% GENERATED by bin/curate/gen_sdem_validation.py from the engine and
% tutorials/steady/membranes/membrane12_sdem_nf270_nacl_traces/constant/experimental/fdl2021_figures.csv
% -- do not edit; regenerate.  Checked for drift by the same script with --check.
\begin{center}\footnotesize
\begin{tabular}{@{}lrrrrr@{}}\toprule
salt & $J_v$ ($\mu$m/s) & Na$^+$ & NH$_4^+$ & Cl$^-$ & NO$_3^-$ \\ \midrule
NaCl & 36.2 & -- & 0.97 & -- & 1.01 \\
MgCl$_2$ & 36.5 & 1.14 & 1.15 & -- & 0.91 \\
Na$_2$SO$_4$ & 33.6 & -- & 0.23 & 1.14 & 1.15 \\
MgSO$_4$ & 34.7 & 0.45 & 0.45 & 1.22 & 1.37 \\
\bottomrule
\end{tabular}
\end{center}

\noindent A ratio of 1 means the trace does not feed back on the field; a
dash means the ion is a dominant ion of that salt.  The paper's own text
names the effect (``trace ions cannot be accounted as authentic traces due
to the very high rejections of SO$_4^{2-}$''); the table is its size.  For a
student it is the lesson of these two cases, and it is why the shipped
MgSO$_4$ witness, run at the paper's concentrations, does not reproduce
the paper's lines while the trace-limit run does.

\subsection*{The two permeances the paper leaves as bounds}

% GENERATED by bin/curate/gen_sdem_validation.py from the engine and
% tutorials/steady/membranes/membrane12_sdem_nf270_nacl_traces/constant/experimental/fdl2021_figures.csv
% -- do not edit; regenerate.  Checked for drift by the same script with --check.
\begin{center}\footnotesize
\begin{tabular}{@{}llrrrrr@{}}\toprule
salt & permeance swept & value ($\mu$m/s) & $J_v$ ($\mu$m/s) & $f$ (Choupo, trace limit) & $f$ (paper's line) & dev.\ (\%) \\ \midrule
NaCl & NH$_4^+$ & 60 & 36.2 & 6.343 & 5.220 & +21.5 \\
NaCl & NH$_4^+$ & 120 & 36.2 & 5.741 & 5.220 & +10.0 \\
NaCl & NH$_4^+$ & 300 & 36.2 & 5.379 & 5.220 & +3.1 \\
NaCl & NH$_4^+$ & 1000 & 36.2 & 5.211 & 5.220 & -0.2 \\
Na$_2$SO$_4$ & Na$^+$ (sets NH$_4^+$) & 0.8 & 27.2 & 64.0 & 72.1 & -11.2 \\
Na$_2$SO$_4$ & Na$^+$ (sets NH$_4^+$) & 3 & 27.2 & 118.4 & 72.1 & +64.3 \\
Na$_2$SO$_4$ & Na$^+$ (sets NH$_4^+$) & 10 & 27.2 & 141.1 & 72.1 & +95.6 \\
Na$_2$SO$_4$ & Na$^+$ (sets NH$_4^+$) & 30 & 27.2 & 148.5 & 72.1 & +106.0 \\
\bottomrule
\end{tabular}
\end{center}

\noindent Under NaCl the NH$_4^+$ line is met when its permeance is of the
order of $1000\,\mu$m/s, an order of magnitude above the stated bound of
60; under Na$_2$SO$_4$ the NH$_4^+$ reciprocal transmission is a strong
function of the Na$^+$ permeance (0.8 $\to$ 30~$\mu$m/s moves it from 64 to
149 at $27\,\mu$m/s), so a value for NH$_4^+$ cannot be read off that
figure without one for Na$^+$.  Neither series is held by the regression
gate, which says so in its own claim.

\subsection*{Points on which we would value the authors' reading}

\begin{enumerate}[nosep]
\item Which curve Fig.~2a shows: Eq.~(1) with $P_s = 6.1\,\mu$m/s does not
      pass through the drawn line or the measured points (previous
      subsection on MgCl$_2$).
\item The right-hand axis of Fig.~3b, read here as 0--18 from its tick
      marks while its labels stop at 8.
\item Whether the lines of the (b) panels were computed from Eq.~(3) (the
      trace limit) or from a full multi-ion solution at $2\times10^{-4}$~M;
      the tables are consistent with the former, and the difference is the
      size of the ``not a trace'' table.
\item The NH$_4^+$ permeance under NaCl: an effective value near
      $1000\,\mu$m/s reproduces the line, the stated bound of 60 does not.
\end{enumerate}

\subsection*{What this study does not establish}

It does not validate the model against the measured \emph{symbols}: the
symbols scatter about the authors' lines by the authors' own fit residual
(up to about $15\,\%$ in $f$ on the last point of Fig.~1a), and a
re-implementation cannot be held to a residual its source accepted.  It
does not test a polarisation film with a field, concentration-dependent
permeances, fixed charge, or any membrane but NF270 at $0.01$~M.  It
inherits the digitisation's reading error, stated above and per point in
the CSV.

\subsection*{Reproducing this section}

The design record of the law and of this comparison is
\file{docs/design/electroneutrality-without-a-new-parameter.md}; the
commands below regenerate every number above.

\Needspace*{8\baselineskip}
\begin{lstlisting}
runCase tutorials/steady/membranes/membrane12_sdem_nf270_nacl_traces   # NaCl witness
runCase tutorials/steady/membranes/membrane13_sdem_nf270_mgso4_traces  # MgSO4 witness
bin/curate/digitise_fdl2021.py --download   # the figures -> the CSV (images stay local)
bin/curate/gen_sdem_validation.py           # every table of this section, from the engine
bin/curate/check_sdem.py                    # the regression gate (arms (h)-(j) hold the lines)
\end{lstlisting}

\subsection*{References}

\begin{enumerate}
\item[{[1]}] A.~Yaroshchuk, M.~L.~Bruening and E.~E.~Lic\'on Bernal,
  ``Solution-diffusion--electro-migration model and its uses for analysis of
  nanofiltration, pressure-retarded osmosis and forward osmosis in multi-ionic
  solutions'', \emph{Journal of Membrane Science} 447 (2013) 463--476,
  doi:10.1016/j.memsci.2013.07.047.  The model Choupo implements as
  \code{transport SDEM;}.
\item[{[2]}] A.~Yaroshchuk and M.~L.~Bruening, ``An analytical solution of
  the solution-diffusion-electromigration equations reproduces trends in ion
  rejections during nanofiltration of mixed electrolytes'', \emph{Journal of
  Membrane Science} 523 (2017) 361--372, doi:10.1016/j.memsci.2016.09.046.
  The trace-ion solution behind the fitted lines of~[3] (its Eq.~(3)).
\item[{[3]}] M.~Fern\'andez de Labastida and A.~Yaroshchuk, ``Nanofiltration
  of Multi-Ion Solutions: Quantitative Control of Concentration Polarization
  and Interpretation by Solution-Diffusion-Electro-Migration Model'',
  \emph{Membranes} 11 (2021) 272, doi:10.3390/membranes11040272, CC~BY~4.0.
  The permeances (Table~1) and the rejection figures this section digitises
  and compares against.
\end{enumerate}
